\section{Prescribing the terminal coupling}
\label{sec:shooting}

The exact transformation of Section~\ref{sec:rg} allows the number of
blocking steps to vary. We now choose the microscopic coupling so that
every trajectory ends at the same coupling $H$. This gives a common
physical normalization and permits comparison of different ultraviolet
cutoffs. The first task is to identify the limiting kinetic increment;
the second is a two-endpoint estimate, using the fixed terminal coupling
at one end and contraction of the remaining data at the other.

\subsection{Initialization and the kinetic increment}

At history zero the unmasked quadratic energy is exactly the
nearest-neighbor energy. We initialize all canonical coefficients and
regular errors at zero. The covering representation is obtained as in
\Rref{rg:trajectories}: join true bad bonds together with the witnesses
for all deleted mask rows, and assign each resulting component the
exponential of minus its energy difference. The difference between the
original energy and the clipped energy is nonnegative. Each true bad
bond contributes at least $cp_j^2$, while every remaining deleted row
is a positive square. The direct-energy reserve therefore pays for the
connected support count and gives the covering cap. The initial bulk
scalar includes the constant converting the chord-square action into
the nearest-neighbor Gibbs weight. This construction uses only chord
bounds and positive squares, and hence applies to $S^2$.

Let $\mathcal C_h$ be the canonical coefficient map constructed in
Section~\ref{subsec:rg-canonical}, where $h$ counts completed blocking
steps. Define its orbit by
\[
 \mathbf P^{(0)}=0,\qquad
 \mathbf P^{(h+1)}=\mathcal C_h(\mathbf P^{(h)}),
\]
and denote the two kinetic increments evaluated on this orbit by
$\alpha_h$ and $\kappa_h$. The map $\mathcal C_h$ is the finite formal
coefficient operation, independent of the precise coupling, regular
errors and covering gas. The exact output representation assigns its
canonical slots to this operation; the discrepancy from the actual
integral remains in the error and gas, with its additional affine
correction assigned to $\Delta$. Thus the canonical slots of an exact
trajectory initialized above follow precisely this orbit. Its actual
coupling recursion is
\begin{equation}
 b_{j-1}=b_j+\alpha_{N-j}+\kappa_{N-j}/b_j+\Delta_j,
 \qquad |\Delta_j|\le C_LH_j^{-1.05},
 \label{eq:shooting-actual-flow}
\end{equation}
which is Equation~\eqref{eq:orig-4} along that trajectory.
The triangular coefficient contraction and the free-history
estimates imply exponential convergence of both increment sequences. The same
limits are obtained if the first step is inclined: after that step the
regular transformations are the same, and their histories converge by
the second comparison in Section~\ref{sec:free}. Denote the limits by
$\alpha_\infty$ and $\kappa_\infty$.

\begin{lemma}[Limiting kinetic increment]
\label{lem:kinetic-increment}
For the nearest-neighbor action and the normalized observations used
here,
\begin{equation}
 \alpha_\infty=-\gamma,
 \qquad \gamma=\frac{\log L}{2\pi}.
 \label{eq:orig-11}
\end{equation}
\end{lemma}

\begin{proof}
We compare a fixed finite-volume Laplace coefficient before and after
blocking. Only the first coefficient is needed, so the second limiting
increment need not be evaluated.

For spins on $S^D$, let $A_D(n)$ be the coefficient of $b^{-1}$ in
$\log Z_b(n,n)$, with $n$ fixed during the expansion. The calculation
\Rref{cal:direct1}, which is stated for general tangent dimension $D$,
gives
\begin{equation}
 4A_D(n)-A_D(2n)
   =-\frac{3D(D-1)}{4\pi}\log n+O(1).
 \label{eq:shooting-bare-coefficient}
\end{equation}
Its action normalization is exactly the nearest-neighbor normalization
of the present model.

Fix a sufficiently large terminal side $m$, put $n=mL^s$, and set
\[
 B_h=\sum_{r<h}\alpha_r,\qquad C_h=\sum_{r<h}\kappa_r,
 \qquad \widehat b_h=b+B_h+C_h/b.
\]
Insert $s$ regular normalized observations with these formal parameters
on the tori of sides $n$ and $2n$. For fixed $s$ they are positive when
$b$ is sufficiently large.
Pin one spin on the final layer. Transitivity of the target sphere
identifies this with integration over the global orientation, up to a
coupling-independent constant. For each fixed $s$ the remaining
integration is finite-dimensional. The aligned saddle has a positive
transverse quadratic form. Outside an aligned neighborhood the action
has a strictly positive excess: small microscopic nearest-neighbor
energy forces alignment, and small observation energy propagates that
alignment through the successive layers. In this neighborhood all
means lie in the normalized-mean branch, so the discontinuous fallback
branch introduces no saddle contribution.

Successive Gaussian completion therefore computes genuine Laplace
coefficients. The computations are precisely the canonical vertex
calculations of Section~\ref{sec:rg}, with cutoffs removed on their
plateaus. The auxiliary normal scalar is integrated with its
normalizing determinant, so each Gaussian pinning factor counts $D$
physical tangent components. Thus the remaining pinning power in the
doubling difference has coefficient $D_0=3D/2$.

For completeness, the uniformity in $s$ needed here is uniformity of
coefficients, not of a bare Laplace remainder. At an output period
$\bar m$, a finite-period coefficient differs from its folded bulk
coefficient only when a complete contraction crosses a period. The
analytic kernel and Wick-coefficient bounds give an exponentially
small error in $\bar m$, with fixed polynomial position factors.
Propagation through one subsequent coefficient calculation costs at
most $C_L$: only finitely many orders occur, and their position sums
are absolutely convergent. The earlier periods are $mL^r$, and
\[
 \sum_{r\ge0} C_L^{r+1}(1+r+mL^r)^C e^{-cmL^r}<\infty.
\]
This is the uniform coefficient estimate in the proof of
\Rref{cal:block1}. It applies here because the projected chart vertices
have the same exponential kernel and finite-contraction bounds, the
canonical slots and increments are bounded, and the last pinned free
operator is uniformly elliptic on the fixed periods.

All extracted bulk scalars cancel in the doubling difference. The
remaining Gaussian power is $D_0\log\widehat b_s$, whose coefficient of
$b^{-1}$ is $D_0B_s$. The bounded final coefficients and the preceding
winding errors consequently give
\[
 4A_D(mL^s)-A_D(2mL^s)
   =D_0\sum_{h<s}\alpha_h+O_{L,m}(1).
\]
Set $D=2$ and compare with
Equation~\eqref{eq:shooting-bare-coefficient}. Division by $s$ and
exponential convergence of $\alpha_h$ give
$\alpha_\infty=-\log L/(2\pi)$.
\end{proof}

\subsection{Shooting to a fixed terminal value}

The remaining layers are indexed by $j=N,N-1,\ldots,0$, so that
$b_N$ is microscopic and $b_0$ is terminal. Put
\[
 c_{\log}=-\frac{\kappa_\infty}{\gamma},
 \qquad
 \vartheta_j=H_j+c_{\log}\log(H_j/H).
\]
The reference sequence satisfies
\begin{equation}
 \vartheta_j-\vartheta_{j-1}
   =\gamma-\frac{\kappa_\infty}{H_j}+O_L(H_j^{-2}).
 \label{eq:shooting-reference}
\end{equation}

\begin{proposition}[Admitted trajectories with prescribed endpoint]
\label{prop:shooting}
After the choices of $L$ and $P_0$, take $H$ sufficiently large. For
every integer $N\ge1$ there is a microscopic coupling $\beta_N>0$
whose $N$ regular transformations remain strictly inside the bands
$[H_j/2,2H_j]$ and end at $b_0=H$. Every such choice satisfies
\begin{equation}
 \beta_N=H+\gamma N+O_{L,H}(\log(2+N)).
 \label{eq:orig-12}
\end{equation}
Along these trajectories,
\begin{equation}
 b_j=\vartheta_j+O_L(1),\qquad 0\le j\le N,
 \label{eq:orig-13}
\end{equation}
uniformly in $N$, with constants bounded as $H$ is increased.
The trajectories starting at the same $\beta_N$ with either inclined
first step are also strictly admitted and satisfy
Equation~\eqref{eq:orig-13}. Their terminal couplings agree with one
another, although they need not equal $H$.
\end{proposition}

\begin{proof}
We give the shooting argument to specify the continuity and uniformity
being used from \Rref{rg:admission}. The exact map is continuous in
the precise input coupling on every finite strictly admitted segment.
For the initial covering representation this follows because its
geometric predicates use the fixed reference thresholds; for each
subsequent step it follows from the uniform product and integration
bounds in Section~\ref{sec:rg}.

On an admitted segment from layer $k$ to layer $j<k$, summing
Equation~\eqref{eq:orig-4}, the bounded second increments, and the
exponentially summable errors in
$\alpha_h+\gamma$ gives
\begin{equation}
 |(b_j-H_j)-(b_k-H_k)|
   \le C_L+C_L(k-j)/H.
 \label{eq:shooting-barrier}
\end{equation}
Vary the initial coupling in $[.6H_N,1.4H_N]$. Call it high if a
strictly admitted initial segment reaches $b_k>1.25H_k$, or if a full
strictly admitted trajectory ends above $H$. Define low with
$b_k<.75H_k$ or a terminal value below $H$. These are nonempty open
subsets of the initial interval: continuity gives openness, and the
two endpoints give nonemptiness.

They are disjoint. An upper and lower mark at layers $j,k$ force a
change at least $(H_j+H_k)/4$ in the quantity $b-H_{\cdot}$. An upper
mark at layer $k$ and a terminal value at most $H$ force a change at
least $H_k/4$; the corresponding lower-mark statement is identical.
These changes contradict Equation~\eqref{eq:shooting-barrier} for
large $H$, because $H_k-H_j=\gamma(k-j)$. Connectedness of the initial
interval therefore supplies an initial value in neither class.
Every step changes the coupling by at most $C_L$. With $H$ sufficiently
large, an exit from an admitted band would have been preceded by one
of the strict internal marks. Thus this value reaches layer zero and
ends exactly at $H$.

To estimate the resulting trajectory, set $D_j=b_j-\vartheta_j$.
The denominator comparison
\[
 |b_j^{-1}-H_j^{-1}|
 \le C\frac{\log(2+H_j/H)+|D_j|}{H_j^2}
\]
and Equation~\eqref{eq:shooting-reference} show that the difference
equation for $D_j$ has summable inhomogeneous terms: exponential
increment errors, $O_L(H_j^{-1.05})$, and
$O_L(H_j^{-2}\log(2+H_j/H))$. Starting from $D_0=0$ and summing
backwards gives, on every finite run,
\[
 \max_j|D_j|
 \le C_L+C_L\left(\sum_{r\ge1}H_r^{-2}\right)\max_j|D_j|.
\]
The coefficient sum is $O_L(H^{-1})$. Increasing $H$ absorbs it and
proves Equation~\eqref{eq:orig-13}, hence
Equation~\eqref{eq:orig-12}.

For an inclined first step, the limiting increments are unchanged
and their transient errors still have bounded sums. Starting from
the same $\beta_N$, the first-exit estimate
\Rref{rg:first-exit-bound} therefore gives
\[
 |D_j|\le C_L+C_L\sum_{r=j+1}^N
 \frac{\log(2+H_r/H)+|D_r|}{H_r^2}
\]
through a hypothetical first exit. Its maximum is bounded by the
same absorption, and the resulting interval lies strictly inside the
admitted band. An exit is impossible. Finally, reflection interchanges
the two inclined prescriptions and preserves every bulk scalar in the
canonical and error normalization. Their running couplings therefore
agree. This scalar equality does not identify their coefficient tensors
in common coordinates; those tensors require the comparison below.
\end{proof}

Fix one such $\beta_N$ for each $N$.  The construction at a prescribed
scale uses any such choice and does not require uniqueness of the shooting
value.  Sections~\ref{sec:trajectories}--\ref{sec:uniqueness} will remove
the choice of trajectory after canonical normalization.

\subsection{Matching the complete endpoint densities}

For two regular trajectories of depths $N\ge K$, align their final
$K$ layers and use the discrepancy variables $u_j,\lambda_j$ of
Section~\ref{sec:rg}. Here $u_j$ measures the normalized coefficient,
error and gas differences, while $\lambda_j$ is the precise coupling
difference. The upper endpoint has bounded shape data; the lower
endpoint satisfies $\lambda_0=0$.

\begin{proposition}[Endpoint matching]
\label{prop:endpoint-matching}
For every sufficiently small fixed $\upsilon>0$, the parameters may be
chosen so that some $\rho<L^{-2+\upsilon}$ satisfies
\begin{equation}
 u_j+|\lambda_j|
 \le C_H(1+K)^C\rho^{K-j},
 \qquad 0\le j\le K,
 \label{eq:orig-14}
\end{equation}
uniformly in $N\ge K$. After the respective bulk scalars have been
removed, write the endpoint densities as $e^{X_i(q)}\Xi_i(q)$.
On every common admitted terminal torus of volume $v$,
\begin{align}
 \norm{X_1-X_2}_\infty&\le D_Hv\delta_K,\nonumber\\
 \sup_x\sum_{\lambda:x\in P_\lambda}
 e^{As_\lambda}\norm{k_{1,\lambda}-k_{2,\lambda}}_\infty
 &\le D_H\delta_K,
 \qquad \delta_K=L^{-(2-\upsilon)K}.
 \label{eq:orig-15}
\end{align}
The constants are independent of depths and periods. The two mirrored
inclined trajectories at depth $N$ have the same endpoint estimates
with $K=N$, when their output lists are expressed in common coordinates.
\end{proposition}

\begin{proof}
This is the two-boundary argument behind \Rref{rg:two-end-bound}, with
the improved contraction and history rate. Let $q$ be the contraction
constant in Equation~\eqref{eq:orig-5}, and put $r=A_0/L^2$. Choose
\[
 \max(q,r)<\rho<L^{-2+\upsilon}
\]
with strict spare width. The coupling coefficient in
Equation~\eqref{eq:orig-5} is bounded by
\[
 \varepsilon_H=\sup_{j\ge0}
 C_L(\log H_j)^C/H_j,
 \qquad \varepsilon_H\longrightarrow0.
\]
Writing $a=(1-\varepsilon_H)^{-1}$, the step inequalities imply
\[
 u_{j-1}\le q u_j+\varepsilon_H|\lambda_j|+f_j,
 \qquad
 |\lambda_j|\le a|\lambda_{j-1}|+aC_Lu_j+af_j,
\]
where $f_j\le C_L\operatorname{poly}(H_j)r^{K-j}$.
Set $w_j=\rho^{K-j}$ and
\[
 U=\max_{0\le j\le K}u_j/w_j,
 \qquad W=\max_{0\le j\le K}|\lambda_j|/w_j.
\]
There is $F_K=C_H(1+K)^C$ such that $f_j\le F_Kw_j$.
Iteration from $K$ for the first inequality and from $0$ for the
second gives
\[
 U\le u_K+\frac{\varepsilon_HW+F_K}{\rho-q},
 \qquad
 W\le\frac{a(C_LU+F_K)}{1-a\rho}.
\]
Take $H$ so that $a\rho<1$ and the product of the two coefficients
coupling $U$ and $W$ is less than $1/2$. Since $u_K$ is bounded by
the common caps, these inequalities prove
Equation~\eqref{eq:orig-14}.

At layer zero the masked slot bounds, the load bounds, and the
free-history discrepancies convert this estimate to the two norms
in Equation~\eqref{eq:orig-15}, exactly as in
\Rref{rg:terminal-difference}. All sup bounds are taken on their
masks; no off-mask extrapolation is used. The factor $(1+K)^C$ is
absorbed into the strict spare width between $\rho$ and
$L^{-2+\upsilon}$. For the inclined comparison, start at layer
$N-1$, where the first-step outputs obey the common caps, and compare
the shared subsequent regular steps. Their terminal couplings agree
by Proposition~\ref{prop:shooting}. The same proof applies, and the
one-step shift is absorbed in the fixed constant.

At every step only a scalar per site has been extracted. The layer-$j$
volume is $vL^{2j}$, so the total scalar is exactly a constant times
$v$, independent of the terminal period. All period-dependent
corrections remain in the winding lists. This proves the claimed
uniformity of the scalar-stripped comparison.
\end{proof}

The endpoint densities are strictly positive: each is obtained by
integrating the positive microscopic density against positive
normalized observations. In standard axis tori their laws also retain
link reflection positivity at every retained block seam. Indeed the
microscopic positive kernel factors across such a seam, and the
observations on its two sides are independent and transform into one
another under reflection. This is the verification following
\Rref{cmp:RG-density}. Translation, quarter-turn, and mirror symmetries
are retained whenever they preserve the period. When tile reflections
are used below, the relevant tile counts are even. No endpoint
reflection positivity is required for oblique periods or after an
inclined first step.
